\chapter{Literature}
\label{ch:literature}

This chapter collects the literature on the unsymmetric finite element method
(UFEM) that the code is built on. All entries are in
\code{literature/references.bib}; the full references with DOIs are listed at the
end of this document. The PDFs are not under version control because they are
publisher-copyrighted or unpublished.

Element acronyms: PP = parametric test / parametric trial (classical
isoparametric); PM = parametric test / metric trial (unsymmetric, ``US-''); MM =
metric / metric; ATF = analytical trial functions.

\section{References followed by the implementation}

\begin{xltabular}{\linewidth}{@{}LP{6.2cm}@{}}
  \toprule
  code & references \\
  \midrule
  unsymmetric (Petrov--Galerkin) element, adjugate form of $\matr B_L$
    (\code{Ke_2D}, \code{Me_2D}, \code{KMe_2D}) &
    \textcite{Rajendran2010technique}; \textcite{Rajendran2003novel} \\
  metric trial functions in corotational skew coordinates
    (\code{ShapeFunctions2D_Metric}, \code{CorotationalSkewCoordinates}) &
    \textcite{Xie2016Modified}; \textcite{Yuan1993New};
    \textcite{Ooi2008Remedies} \\
  Lee \& Bathe patch tests (\code{code/InputFilesAbaqus},
    \code{SpatialLoadingFunction}) &
    \textcite{Lee1993effects}; \textcite{Rajendran2003novel} \\
  adaptive hybrid explicit method (\code{AdaptiveHybridExplicitMethod}) &
    \textcite{Soares2020novel} \\
  stability of time integration with unsymmetric matrices &
    \textcite{Wang2022Temporal}; \textcite{Zaehringer2026On} \\
  \bottomrule
\end{xltabular}

\section{Annotated bibliography}

% Two-column table: reference and key content
\newcommand{\littable}[1]{%
  \begin{xltabular}{\linewidth}{@{}P{4.4cm}L@{}}
    \toprule
    reference & key content \\
    \midrule
    \endhead
    #1
    \bottomrule
  \end{xltabular}}

\subsection{Axisymmetric UFEM}
\littable{
  \textcite{Zaehringer2024finite} & 8-node PG-Q8: serendipity test functions, metric
    trial functions in skew coordinates. Exact for quadratic fields on distorted
    meshes; for the thick-walled cylinder (non-polynomial $u_r$) not generally
    better than Q8. \\
  \textcite{Dou2025novel} & UQ4-RPIM; extension to axisymmetry with a modified
    polynomial basis for hollow bodies. \\
  \textcite{Huang2020incompatible} & IUQ4; axisymmetric extension with incompatible
    test modes that pass the axisymmetric patch test. \\
}

\subsection{Foundations of UFEM and Petrov--Galerkin elements}
\littable{
  \textcite{Rajendran2003novel} & US-QUAD8, the seminal paper: serendipity
    (compatibility-enforcing) test functions with metric (completeness-enforcing)
    trial functions; exact quadratic fields under angular, curved-edge and
    mid-node distortion. \\
  \textcite{Rajendran2000completeness} & higher-order completeness conditions;
    why distorted Q8 elements lose quadratic completeness. \\
  \textcite{Rajendran2010technique} & three sufficient conditions for distortion
    immunity (compatible test, complete trial, no $\det\matr J$ in the
    denominator); PP/PM/MM naming. \\
  \textcite{Ooi2006unsymmetric} & comprehensive study of US-QUAD8/US-HEXA20:
    strong-form patch tests, geometric nonlinearity, frame dependence, P-matrix
    singularity, complex eigenvalues. \\
  \textcite{Prathap2007The} & the PM formulation is not variationally correct but
    recovers stresses as a best fit in metric space. \\
  \textcite{Cen2019Some} & review of hybrid stress/displacement-function elements
    and the improved UFEM (US-ATF family). \\
  \textcite{Chen2021implementierung} & Abaqus UEL implementation of PM-QUAD8 and
    PM-HEXA20; cantilever and Kirsch-plate benchmarks under distortion. \\
}

\subsection{Skew coordinates and frame invariance}
\littable{
  \textcite{Xie2016Modified} & UQ8*, a frame-invariant UQ8 with metric trial
    functions in non-dimensional skew coordinates; Trefftz-UFEM TQ4/TH8. \\
  \textcite{Yuan1993New} & origin of the skew (oblique) coordinates defined at the
    element centre. \\
  \textcite{Ooi2008Remedies} & US-QUAD8 is frame-dependent for the cubic terms and
    its P-matrix can become singular; remedies. \\
  \textcite{Cowan2014Rotationally} & metric trial functions in a local convected
    frame (rotational invariance); finite-deformation elasto-plasticity. \\
  \textcite{Wisniewski2008improved} & Hellinger--Reissner 4-node elements with
    stresses in skew coordinates. \\
}

\subsection{Plane elements}
\littable{
  \textcite{Rajendran2004mesh} & comparison of PP, MM, PM and MP 8-node elements. \\
  \textcite{Liew2006quadratic} & US-TRIA6; free vibration with an unsymmetric mass
    matrix. \\
  \textcite{Cen2012shape} & US-ATFQ8 with Airy-function analytical trial functions. \\
  \textcite{Cen2015unsymmetric} & US-ATFQ4, ``breaking through'' MacNeal's
    theorem. \\
  \textcite{Shang2018quadrilateral} & US-Q4 with drilling DOFs and stress trial
    functions. \\
  \textcite{Shang2018airy} & US-Q8 for orthotropic problems (Airy stress
    solutions). \\
  \textcite{Huang2022Unsymmetric} & UQ6, UH11: Wilson incompatible modes, explicit
    trial functions. \\
  \textcite{Yang2023Two} & UQ10, UQ10/E4: cubic completeness through incompatible
    modes / EAS. \\
}

\subsection{Solid and solid-shell elements}
\littable{
  \textcite{Ooi2004hexahedron} & US-HEXA20. \\
  \textcite{Zhou2017unsymmetric} & US-ATFH8. \\
  \textcite{Huang2018unsymmetric} & US-ATFHS8 solid-shell. \\
  \textcite{Wu2023unsymmetric} & UHS11 solid-shell. \\
  \textcite{He2019generalized} & generalized modal element method (US-MEM8S). \\
  \textcite{Pfefferkorn2021mesh} & Petrov--Galerkin EAS (Q1U/E4, H1U/E12);
    MacNeal's theorem revisited. \\
}

\subsection{Nonlinear extensions}
\littable{
  \textcite{Ooi2007extension} & US-QUAD8/US-HEXA20, total Lagrangian. \\
  \textcite{Li2018high} & US-ATFQ4, updated Lagrangian. \\
  \textcite{Li2019unsymmetric} & US-ATFHS8, corotational. \\
  \textcite{Li2020hyperelastic} & US-ATFQ4, hyperelastic. \\
  \textcite{Chen2021unsymmetric} & US-QUAD8, elasto-plastic. \\
  \textcite{Pfefferkorn2023hourglassing} & PG-EAS H1U/h12, large deformations. \\
  \textcite{Wang2023high} & US-ATFH8 for soft tissues. \\
  \textcite{Wu2026unsymmetric} & UHS11, hyperelastic. \\
}

\subsection{Dynamics and time integration}
\littable{
  \textcite{Wang2022Temporal} & with non-symmetric $\matr M$/$\matr K$ energy is not
    conserved and Newmark can be unconditionally unstable; energy-conserving
    scheme. \\
  \textcite{Zaehringer2026On} & UQ8* in elastodynamics: standard Newmark is
    unstable, Wang--Hillman can drift, a mixed velocity--strain formulation
    conserves energy and stays distortion-insensitive. \\
  \textcite{Kumar2011Use} & transient (impact) application of PM20 with Newmark. \\
  \textcite{Soares2020novel} & adaptive hybrid explicit method (method 109). \\
}

\subsection{Plates, beams and shells}
\littable{
  \textcite{Ju2025unsymmetric} & UQP4-RPIM Reissner--Mindlin plate. \\
  \textcite{Shang2022Trefftz} & US-TP4 Trefftz plate element. \\
}

\subsection{Couple-stress and flexoelectric continua}
\littable{
  \textcite{Shang2019simple} & US-Q4-CS. \\
  \textcite{Shang2020hexahedral} & 8-node hexahedron (modified couple stress). \\
  \textcite{Shang2022couple} & US-Q8, US-Q8L. \\
  \textcite{Wu2023Penalty} & penalty $C^0$ Q8/H20 (consistent couple stress). \\
  \textcite{Deng2025corotational} & corotational flexoelectric membrane. \\
  \textcite{Shang2025Mesh} & Trefftz Kirchhoff plate (consistent couple stress). \\
}

\subsection{Mesh distortion: theory and assessment}
\littable{
  \textcite{Lee1993effects} & angular distortion destroys the quadratic
    completeness of serendipity elements. \\
  \textcite{MacNeal1992On} & MacNeal's theorem: a symmetric element must either
    fail the patch test or lock. \\
  \textcite{Rajendran2007Mesh} & strong-form patch tests; only PM-QUAD8 passes all. \\
  \textcite{Prathap2006mesh} & best-fit paradigm explanation of the distortion
    immunity of PM elements. \\
  \textcite{Mukherjee2011improved} & variationally correct, distortion-immune
    3-node bar. \\
  \textcite{Manju2019function} & function-space analysis of the distorted 3-node
    Timoshenko metric beam. \\
  \textcite{Engvall2020Mesh} & mesh quality metrics for high-order
    (Bernstein--Bézier) elements. \\
  \textcite{Quiriny2024The} & tempered FEM: a lower bound on $\det\matr J$ instead of
    a change of test/trial spaces. \\
  \textcite{Quiriny2026Taming} & tempered FEM for sliver tetrahedra. \\
}

\section{Synthesis}

\begin{enumerate}
  \item \textbf{Why UFEM.} Conventional (Bubnov--Galerkin) serendipity elements lose
    their quadratic completeness under angular and curved-edge distortion
    \parencite{Lee1993effects}. MacNeal's theorem shows that no symmetric element
    can both pass the patch test and be distortion-insensitive in bending
    \parencite{MacNeal1992On}. Petrov--Galerkin elements avoid this by combining
    compatible (parametric) test functions with Cartesian-complete (metric) trial
    functions \parencite{Rajendran2003novel,Rajendran2010technique}.
  \item \textbf{Sufficient conditions for distortion immunity}
    \parencite{Rajendran2010technique}: compatible test functions, trial functions
    complete in Cartesian coordinates, no $\det\matr J$ in the denominator of the
    integrand. This code meets all three for arbitrary $p$: the trial functions are
    tensor-product Lagrange polynomials in skew coordinates, which contain all
    polynomials of total degree $\le p$ in $(x, y)$, and $\matr B_L$ is built with
    $\adj\matr J$. Any displacement field of total degree $\le p$ is reproduced
    exactly on arbitrarily distorted meshes (\code{tests/integration/PatchTest.m}).
  \item \textbf{Frame invariance.} Metric trial functions in global coordinates are
    frame-dependent and their Vandermonde matrix can become singular
    \parencite{Ooi2008Remedies}. Skew coordinates
    \parencite{Yuan1993New,Xie2016Modified} fix the frame dependence; they are the
    default here. The Vandermonde matrix can still be ill-conditioned for
    degenerate elements, e.g.\ three nearly collinear corners.
  \item \textbf{Variational inconsistency.} The PM formulation is not variationally
    correct \parencite{Prathap2007The}: the stiffness matrix is unsymmetric and the
    energy optimality of Bubnov--Galerkin is lost. For fields outside the trial
    space the unsymmetric element can be slightly less accurate than the classical
    one on coarse meshes \parencite{Zaehringer2024finite}, but it keeps the optimal
    convergence rate on distorted meshes.
  \item \textbf{Dynamics.} Unsymmetric $\matr K$ (and $\matr M$) break the
    equivalence of energy and momentum conservation; standard Newmark can become
    unstable and generalized eigenvalues can be complex
    \parencite{Wang2022Temporal,Ooi2006unsymmetric,Zaehringer2026On}. With GLL nodes
    and nodal quadrature the unsymmetric mass matrix is nevertheless diagonal,
    because both the parametric and the metric functions are interpolatory at the
    quadrature points.
  \item \textbf{Open questions.} Element-level stability (conditioning of the metric
    Vandermonde matrix, discrete inf-sup constant) on strongly distorted high-order
    elements; energy-consistent time integration for unsymmetric spectral elements;
    a convergence theory for these variationally inconsistent schemes; extension to
    triangles (\code{ShapeFunctions2D_Metric_Tria} is an unfinished start), 3D and
    axisymmetry.
\end{enumerate}
